\section{GPU implementation}
\label{sec:gpu_impl}

\subsection{Parallelisation strategy}

The interaction-picture evolution
\begin{equation}
  \frac{d\rho}{dt} \;=\; i[\rho,\, H_\mathrm{I}(t)]
                       \;-\; \{\Gamma(t),\, \rho\}
                       \;+\; \mathcal{F}_\mathrm{int}(t)
  \label{eq:vN}
\end{equation}
is independent across zenith trajectories. Each trajectory is mapped
to a single CUDA thread block, with the energy and flavor dimensions
striped across threads within the block. State -- the SU($N_\nu$)
density-matrix coefficients $\rho^c$ for each $(E_i, \rho)$ -- lives
in shared memory. Profile information (matter density, electron
fraction $Y_e$, target species number densities) is uploaded as
Akima-spline coefficients in a per-block constant-memory buffer.

% TODO: schematic figure of the block layout.

\subsection{Adaptive step control with substage cascade refresh}

The CPU \nusquids\ uses RKF45 from \textsc{GSL}, which evaluates the
right-hand side of Eq.~\eqref{eq:vN} multiple times per accepted step
and -- crucially -- recomputes the cascade-source factors at every
substage by invoking \texttt{PreDerive(t)}. On the GPU we adopt a
Richardson step-doubling scheme: a full step of size $h$ is computed
in parallel with two half-steps of size $h/2$, and the difference
between the two estimates serves as the truncation-error proxy
that drives accept/reject and step-size adaptation.

To match the CPU's per-substage refresh behaviour, the GPU integrator
performs the following sequence per macro step:
\begin{enumerate}
  \item Take a predictor half-step using start-of-step cascade
    factors. Store the predicted state at $t+h/2$ in shared memory.
  \item Cooperatively refresh the cascade source terms (NC, tau
    decay, Glashow) at $t+h/2$ using the predicted mid-state. The
    refresh kernels evaluate the energy-coupled reductions across
    threads.
  \item Recompute both Richardson branches -- the full step $s_f$ and
    the two half-steps culminating in $s_h$ -- using the refreshed
    cascade factors. This guarantees that $s_f$ and $s_h$ solve the
    same forced ODE, which is the prerequisite for the Richardson
    extrapolation to be unbiased.
  \item Form the corrected state $s_h + (s_h - s_f)/15$ and the local
    error estimate; reduce across the block to decide accept/reject.
\end{enumerate}
Substep-by-substep refreshes (4 per RK4 step) further tighten the
agreement; we adopt this as the default. Without per-substage
refresh, tau regeneration agreement degrades from sub-percent to
several tens of percent because of the two-stage nature of the source
(charged-current production followed by decay-spectrum redistribution).

\subsection{Cascade-source kernels}

The three cascade contributions are computed by separate cooperative
kernels: \texttt{computeNCCascadeSU3} (one-stage, $\mathcal{O}(n_E^2)$
per refresh), \texttt{computeTauRegenSU3} (two-stage, $\mathcal{O}(n_E^3)$
production-then-decay), and \texttt{computeGlashowCascadeSU3} (sharp
resonance at $\sim 6.3\,\mathrm{PeV}$, $\mathcal{O}(n_E^2)$). Profile
quantities -- target number densities, $Y_e$, and the Glashow electron
number density $n_e$ -- are pre-evaluated once per substage by a
single thread and broadcast through shared memory, eliminating per-thread
Akima spline evaluations.

% TODO: short subsection on the Glashow num_e body-composition path
% (include the full CPU-mirroring formula num_e = density / (m_e + m_p + m_n*(1-y_e)/y_e))

\subsection{CPU correctness as ground truth}

Throughout the GPU port we treat the CPU \nusquids\ implementation as
the oracle. All physics constants are sourced from the SQuIDS
\texttt{Const} natural-units system; no conversion factors are
hardcoded on the device. Where the GPU requires a precomputed quantity
that the CPU computes inline (e.g.\ the Glashow $n_e$ for bodies with
non-trivial composition), the value is computed CPU-side using the
same formula and uploaded as profile data, never recomputed
independently on the device.
